\documentclass[10pt]{article}
\usepackage[margin=1.4cm]{geometry}
\usepackage{booktabs}
\usepackage{graphicx}
\usepackage[table]{xcolor}

\definecolor{lorarow}{RGB}{244,247,251}

\title{Reconstruction Interventions by Editing Method}
\date{}

\begin{document}
\maketitle

This report records the editor-specific interventions considered after the
LoRA+DirectInv models improved CLIP editing scores but remained behind
DirectInv on reconstruction. The purpose was not to change LoRA training, but
to test whether each editor consumes the LoRA-inverted trajectory in a way that
is mismatched with its original DirectInv implementation.

All numerical results use SD1.5 and the same 700-image PIE-Bench evaluation.
Values are mean $\pm$ standard deviation. Structure distance and LPIPS are
scaled by $10^3$, MSE by $10^4$, and SSIM by $10^2$. Bold denotes the better
value in each matched DirectInv/LoRA+DirectInv pair after rounding; arrows give
the preferred metric direction. Structure uses all 700 images. PSNR, LPIPS,
MSE, and SSIM use the 556 images with a non-empty unedited mask.

\section*{Interventions and outcomes}

\begin{table}[htbp]
  \centering
  \small
  \setlength{\tabcolsep}{4pt}
  \begin{tabular}{p{1.7cm}p{4.2cm}p{5.1cm}p{3.4cm}}
    \toprule
    Editor & Suspected mismatch & Intervention & Outcome \\
    \midrule
    P2P & No new editor mismatch was isolated; cycle-loss training had already
    produced the strongest reconstruction result. & No additional editor-side
    change in this round. & Not tested here. \\
    \addlinespace
    MasaCtrl & Standard DirectInv inverts its source branch with the empty
    prompt, while fused LoRA inversion uses the source prompt. & Verify both
    conditioning paths; next test LoRA inversion with its unconditional adapter
    while leaving MasaCtrl generation unchanged. & Diagnosis confirmed;
    intervention not run. \\
    \addlinespace
    PnP & The injected reference latent used an empty text embedding despite
    originating from source-conditioned inversion. & Condition only PnP's
    injected reference branch on the source prompt; preserve target CFG and
    injection schedules. & Did not improve reconstruction; CLIP improved. \\
    \addlinespace
    Pix2Pix-Zero & A raw sum over cross-attention errors can let large maps or
    numerous layers dominate the latent update. & Average each layer's MSE
    equally, cap the latent-gradient norm at 1.0, and sweep attention guidance.
    & Failed: reconstruction and editing degraded. \\
    \bottomrule
  \end{tabular}
  \caption{Editor-specific reconstruction hypotheses and their status.}
  \label{tab:intervention-summary}
\end{table}

\section*{PnP: source-conditioned reference features}

At each denoising step, PnP evaluates a three-branch batch: the saved reference
latent, the target unconditional branch, and the target conditional branch.
Self-attention keys/queries and convolutional features from the first branch
are injected into the target branches. The tested change replaced the empty
text embedding of only that first branch with the source-prompt embedding. It
did not alter the target prompt, target CFG, noisy trajectory, or injection
timesteps. Partially projected trajectories were deliberately not tested.

\begin{table}[htbp]
  \centering
  \scriptsize
  \setlength{\tabcolsep}{2.5pt}
  \caption{PnP with its injected source/reference branch conditioned on the
  source prompt instead of the empty prompt.}
  \label{tab:pnp-source-reference}
  \resizebox{\textwidth}{!}{%
  \begin{tabular}{llrrrrrrr}
    \toprule
    Editor & Inversion & Structure $\downarrow$ & PSNR $\uparrow$ & LPIPS $\downarrow$ & MSE $\downarrow$ & SSIM $\uparrow$ & CLIP target $\uparrow$ & CLIP edit $\uparrow$ \\
    \midrule
    PnP & DirectInv
      & $\mathbf{23.38 \pm 19.79}$ & $\mathbf{22.21 \pm 4.31}$
      & $\mathbf{108.38 \pm 61.33}$ & $\mathbf{83.84 \pm 59.81}$
      & $\mathbf{80.21 \pm 12.80}$ & $25.75 \pm 3.82$ & $22.71 \pm 4.93$ \\
    \rowcolor{lorarow}
    PnP & LoRA+DirectInv
      & $24.23 \pm 19.50$ & $22.13 \pm 4.55$
      & $112.39 \pm 63.94$ & $85.93 \pm 60.86$
      & $79.85 \pm 12.86$ & $\mathbf{25.91 \pm 3.81}$
      & $\mathbf{22.90 \pm 4.95}$ \\
    \bottomrule
  \end{tabular}%
  }
\end{table}

\paragraph{Verdict: no reconstruction improvement.}
Within the new source-conditioned setup, LoRA+DirectInv is worse than DirectInv
on all five reconstruction metrics, while CLIP target rises from 25.75 to 25.91
and CLIP edit from 22.71 to 22.90. Compared with the previous empty-reference
LoRA+DirectInv result, Structure improves slightly (24.38 to 24.23), but PSNR,
LPIPS, MSE, and SSIM worsen (22.36 to 22.13, 110.28 to 112.39, 81.87 to 85.93,
and 79.88 to 79.85). The corresponding DirectInv control also fails to improve
reconstruction relative to its empty-reference result. The change improves
semantic edit alignment, not preservation.

\section*{Pix2Pix-Zero: normalized attention guidance}

Pix2Pix-Zero first stores cross-attention maps, then optimizes the current
latent during editing to match those maps. The legacy objective sums squared
errors over all elements of all layers. The tested objective instead computes
mean-squared error within each layer and averages layers equally; it also caps
the latent-gradient norm at 1.0. Attention-guidance amounts 0.025, 0.050, 0.075,
and 0.100 were applied identically to DirectInv and LoRA+DirectInv. This uses
Pix2Pix-Zero's existing reference maps and does not add another inversion pass.

\begin{table}[htbp]
  \centering
  \scriptsize
  \setlength{\tabcolsep}{2.5pt}
  \caption{Pix2Pix-Zero with equal per-layer attention-loss weighting and a
  latent-gradient norm cap of 1.0. The guidance column is the Pix2Pix-Zero
  attention-guidance amount, not diffusion CFG.}
  \label{tab:pix2pix-normalized}
  \resizebox{\textwidth}{!}{%
  \begin{tabular}{llrrrrrrr}
    \toprule
    Guidance & Inversion & Structure $\downarrow$ & PSNR $\uparrow$ & LPIPS $\downarrow$ & MSE $\downarrow$ & SSIM $\uparrow$ & CLIP target $\uparrow$ & CLIP edit $\uparrow$ \\
    \midrule
    0.025 & DirectInv
      & $\mathbf{79.73 \pm 49.91}$ & $\mathbf{18.44 \pm 4.39}$
      & $\mathbf{223.36 \pm 120.66}$ & $\mathbf{201.66 \pm 151.16}$
      & $\mathbf{70.95 \pm 16.26}$ & $\mathbf{21.96 \pm 4.61}$
      & $20.68 \pm 4.72$ \\
    \rowcolor{lorarow}
    0.025 & LoRA+DirectInv
      & $81.87 \pm 50.08$ & $18.17 \pm 4.23$
      & $229.76 \pm 122.37$ & $211.61 \pm 157.69$
      & $70.35 \pm 16.28$ & $\mathbf{21.96 \pm 4.78}$
      & $\mathbf{20.82 \pm 4.68}$ \\
    \addlinespace
    0.050 & DirectInv
      & $\mathbf{79.73 \pm 49.91}$ & $\mathbf{18.44 \pm 4.39}$
      & $\mathbf{223.36 \pm 120.66}$ & $\mathbf{201.66 \pm 151.16}$
      & $\mathbf{70.95 \pm 16.26}$ & $\mathbf{21.96 \pm 4.61}$
      & $20.68 \pm 4.72$ \\
    \rowcolor{lorarow}
    0.050 & LoRA+DirectInv
      & $81.87 \pm 50.08$ & $18.17 \pm 4.23$
      & $229.76 \pm 122.37$ & $211.61 \pm 157.69$
      & $70.35 \pm 16.28$ & $\mathbf{21.96 \pm 4.78}$
      & $\mathbf{20.82 \pm 4.68}$ \\
    \addlinespace
    0.075 & DirectInv
      & $\mathbf{79.92 \pm 50.41}$ & $\mathbf{18.45 \pm 4.43}$
      & $\mathbf{223.48 \pm 120.68}$ & $\mathbf{201.74 \pm 150.73}$
      & $\mathbf{70.96 \pm 16.24}$ & $\mathbf{21.93 \pm 4.62}$
      & $20.69 \pm 4.74$ \\
    \rowcolor{lorarow}
    0.075 & LoRA+DirectInv
      & $81.90 \pm 50.01$ & $18.17 \pm 4.27$
      & $229.59 \pm 122.45$ & $211.71 \pm 157.57$
      & $70.35 \pm 16.28$ & $\mathbf{21.93 \pm 4.81}$
      & $\mathbf{20.80 \pm 4.64}$ \\
    \addlinespace
    0.100 & DirectInv
      & $\mathbf{79.93 \pm 50.41}$ & $\mathbf{18.45 \pm 4.43}$
      & $\mathbf{223.48 \pm 120.68}$ & $\mathbf{201.74 \pm 150.73}$
      & $\mathbf{70.96 \pm 16.24}$ & $\mathbf{21.93 \pm 4.62}$
      & $20.69 \pm 4.74$ \\
    \rowcolor{lorarow}
    0.100 & LoRA+DirectInv
      & $81.90 \pm 50.01$ & $18.17 \pm 4.27$
      & $229.59 \pm 122.45$ & $211.71 \pm 157.57$
      & $70.35 \pm 16.28$ & $\mathbf{21.93 \pm 4.81}$
      & $\mathbf{20.80 \pm 4.64}$ \\
    \bottomrule
  \end{tabular}%
  }
\end{table}

\paragraph{Verdict: failed.}
The four guidance values are almost indistinguishable: 0.025 and 0.050 round to
the same results, as do 0.075 and 0.100. More importantly, even normalized
DirectInv is much worse than the reproduced legacy SD1.5 DirectInv control:
Structure changes from 49.09 to at least 79.73, PSNR from 21.57 to at most
18.45, LPIPS from 138.39 to at least 223.36, MSE from 129.77 to at least 201.66,
and SSIM from 77.73 to at most 70.96. CLIP target and edit also fall from 23.47
and 21.08 to at most 21.96 and 20.69. LoRA+DirectInv remains worse on all
reconstruction metrics and recovers only a small edited-region CLIP advantage
within each normalized pair.

Because layer normalization, gradient clipping, and guidance recalibration were
changed together, this experiment rejects their current combination rather
than identifying one component as the sole cause. The very small guidance
sensitivity suggests that converting a summed loss to a mean reduced the update
scale too strongly. A future ablation should log the pre/post-clipping gradient
norms and calibrate the normalized update to the legacy update magnitude before
changing its direction.

\section*{MasaCtrl: conditioning mismatch identified, not yet tested}

The reproduced standard DirectInv+MasaCtrl control already uses empty-prompt
inversion: MasaCtrl constructs the pair \texttt{[empty, target]}, and
DirectInv selects the first embedding for its inversion trajectory. Therefore
rerunning the base control with an empty prompt would duplicate an existing
result.

The fused path is different: its LoRA trajectory is generated with the source
prompt, after which DirectInv offsets and MasaCtrl generation use the standard
empty/target contexts. This means the current comparison changes both the
inversion model and its conditioning. The clean proposed experiment is
base+empty versus LoRA+empty, using the unconditional adapter during LoRA
inversion while leaving offset calculation and MasaCtrl generation unchanged.
That experiment was not part of these jobs, so there is no positive or negative
MasaCtrl result to report yet.

\section*{P2P}

No new P2P editor-side intervention was evaluated in this round. P2P already
benefited most from the cycle-loss LoRA, so the investigation prioritized the
editors whose reconstruction gap remained largest. Its status in
Table~\ref{tab:intervention-summary} is a scope decision, not evidence that no
further P2P improvement is possible.

\section*{Source-image traits of worst reconstructions}

This descriptive analysis uses the current cycle-loss r16 uniform
LoRA+DirectInv outputs. For each editor, the 25 images that most degrade the
five reconstruction metrics jointly were compared with the other 531 PIE-Bench
images having all five valid metrics. No edit-category labels are used here:
these are source-image properties only. Signed differences are standardized
mean differences (worst minus remaining); they are descriptive, not causal.

\begin{table}[htbp]
  \centering
  \small
  \setlength{\tabcolsep}{4pt}
  \begin{tabular}{lrrrr}
    \toprule
    Editor & Artificial & Luminance contrast & Edge detail & Unedited area \\
    & worst / remaining & $z$ & $z$ & worst / remaining \\
    \midrule
    Pix2Pix-Zero & 28\% / 39\% & $+0.10$ & $+0.41$ & 85\% / 67\% \\
    MasaCtrl     & 52\% / 38\% & $+0.75$ & $+0.53$ & 84\% / 67\% \\
    PnP          & 68\% / 37\% & $+0.16$ & $+0.78$ & 83\% / 67\% \\
    \bottomrule
  \end{tabular}
  \caption{Source-image characteristics of the worst reconstruction cases.
  Higher edge detail is the only common visual signal. The larger unedited area
  shows these are not failures caused by evaluating a tiny residual region.}
  \label{tab:worst-source-traits}
\end{table}

Pix2Pix-Zero's worst cases skew toward natural images and moderately higher
edge detail. MasaCtrl's worst cases are more often artificial and have markedly
higher luminance contrast. PnP has the strongest artificial-graphics and
edge-detail enrichment. Mean saturation and grayscale entropy showed no
consistent cross-editor effect.

\section*{Overall conclusion}

Neither completed editor-side modification solved reconstruction. PnP source
conditioning shifted the trade-off toward better semantic editing but slightly
worse preservation. The current Pix2Pix-Zero normalization-and-clipping setup
caused a large regression and should not replace the legacy objective. The
MasaCtrl conditioning hypothesis remains the most direct untested control; it
must be evaluated with matched empty conditioning before attributing its gap to
the LoRA itself.

\end{document}
